\scanpage{221}
\mspage{137}
% archive sheet number: 220
% element: text | top, continuation | resulte aussitot du
résulte aussitôt du

% element: text | top, corollary heading and rewritten hypothesis | Corollaire 4.5.16
\uline{Corollaire} 4.5.16 (de 4.5.15). Dans les conditions de 4.5.1\deleted{\unsure{3}}5, \deleted{(\unsure{où|ou} \unsure{l'hypothèse|la condition} C) de 4.5.13 \unsure{(3 C)|3 C})} \inserted{(i.e. la condition C) de 4.5.13.3°)},
% element: deletion | upper third, crossed line with dotted strokes beneath | soit x un objet de M
% Short separate strokes run below this crossed line and below existe; unsure alternatives explicitly preserve cancelled versus reinstated state.
\unsure{\deleted{soit $x$ un objet de $M$, et supposons qu'il existe}|soit $x$ un objet de $M$, et supposons qu'il existe} une famille \uline{épimorphique} \uuline{universelle} $(x_\alpha\longrightarrow x)_{\alpha\in A}$ de flèches \inserted{de $M$}, avec $x_\alpha\in\mathrm{Ob}\,\mathrm{Filt}'{}^\pi$ pour $\forall\alpha\in A$, et $\operatorname{card}A\leq\pi$. Alors $x\in\mathrm{Ob}\,\mathrm{Filt}'{}^\pi$.

% element: text | middle upper, proof | En effet l hyp sur x alpha
En effet, l'hyp. \unsure{sur|\deleted{sur}} $x_\alpha$ \deleted{implique} \unsure{équivaut|équivaut à}, en vertu de 4.5.15, qu'il existe une famille \inserted{épim. stricte universelle} \unsure{de} flèches
% element: diagram | middle, covering family of x alpha | x alpha beta indexed by beta in B alpha to x alpha
\diagram{\begin{tikzcd}[ampersand replacement=\&]
(x_{\alpha\beta})_{\beta\in B_\alpha}\arrow[r]\&x_\alpha,
\end{tikzcd}}
% element: text | middle, cardinal bound | avec card B alpha
\deleted{\unsure{pour un}} \deleted{de} \inserted{avec} $\operatorname{card}B_\alpha\leq\pi$ et les $x_{\alpha\beta}$ dans $C$. Soit
% element: display | middle, disjoint union of index sets | B equals coproduct B alpha
\[
B=\coprod_{\alpha\in A}B_\alpha,
\]
% element: text | middle lower, composite family | on a alors card B
on a alors $\operatorname{card}D\leq\pi^2=\unsure{\pi}$, et d'autre part la famille des flèches composées
% element: diagram | lower middle, composite arrows | x alpha beta to x alpha to x
\diagram{\begin{tikzcd}[ampersand replacement=\&]
x_{\alpha\beta}\arrow[r]\&x_\alpha\arrow[r,dash]\&x
\end{tikzcd}}
% element: text | lower third, conclusion | indexee par B
indexée par $B$ est épimorphique stricte universelle (par transitivité de ce type de familles), donc $x\in\mathrm{Ob}\,\mathrm{Filt}'{}^\pi$, qed.

% element: drawing | space above final paragraph | isolated oblique stroke
\figure{\begin{tikzpicture}\draw (0,0)--(0.35,-0.5);\end{tikzpicture}}

% element: text | foot, transition to characterization | Avec tout ca
Avec tout ça, on n'a toujours pas donné de caractérisation de $M_\pi$ lui-même, qui est \unsure{meilleur|meilleure} que $\mathrm{Filt}^\pi$ et $\mathrm{Filt}'{}^\pi$ du point de vue \unsure{catégoriques|catégorique}, \unsure{car} \inserted{seuls} les $M_\pi$ \deleted{\unsure{forment}} définissent une filtration cardinale, et de plus ils ont les propriétés de stabilité désirables dans \unsure{ces|les}
